\documentclass[12pt]{article}
\sloppy
\hfuzz=10pt
\tolerance=1000
\usepackage{amsmath, amssymb}
\usepackage{tikz-cd}
\usepackage{relsize}
\usepackage{amsthm}

\theoremstyle{definition}

\newtheorem{definition}{Definition}[section]
\newtheorem{lemma}[definition]{Lemma}
\newtheorem{theorem}[definition]{Theorem}
\newtheorem{example}{Example}
\newtheorem*{lemmastar}{Lemma}

\title{On Algebraic Integral For Dummies}
\author{Nikita Umerov}
\date{April 17, 2026}

\begin{document}

\maketitle

\begin{abstract}
We explain the essence of the integral axiom and provide examples of integrating elementary 
functions using the new integration method.
\end{abstract}

\section*{Introduction}

My paper "On Algebraic Integral" \cite{umerov2025} has not received the circulation it deserves. 
Moreover, it appears that guys from the so-called "academic community" — an international 
governmental organization designed to slow down scientific progress to protect 
the governments from dissent — these guys have failed to understand my discovery. 
At the very least, I have received no significant response from them.

\vspace{5pt}

That is why I decided to write a simplified version of the paper — a version for dummies — 
to show the essence of the discovery in the most basic terms. 
I also want to highlight some of its possible applications. 
Of course, the applications of that discovery are not limited to the facts presented in the paper. 
Many of them I cannot even imagine — just as Arthur Cayley could not have imagined the use of matrices 
in artificial intelligence. I will only point out the most obvious ones.

\vspace{5pt}

Moreover, as a researcher, I see my primary mission as exposing the bankruptcy of the academic 
community (together with the contemporary educational system) and its immeasurable harm to society. 
I am currently working on a book on this very subject and do not plan to be 
greatly distracted from that work. I am writing this article only as an explanation 
to accompany the main paper "On Algebraic Integral" \cite{umerov2025} — in response to requests I 
have received from several people. 

\vspace{5pt}
\newpage

\section*{What is the integral axiom?}

In the first part of the paper, we will show that the integral is not a function of a single variable, 
but an operation $f \triangleright g = \int f \, dg$. While the method of change of variables is 
nothing other than the distributivity of composition over that operation:

\[
(f \triangleright g) \circ h = (f \circ h) \triangleright (g \circ h)
\]

\noindent \textbf{Let's begin.}

\vspace{20pt}
\noindent Let $\langle R, \cdot, \circ \rangle$ be a near-ring. That is, a ring with only right distributivity. 
That is, $R$ such a structure that 

\[
\forall f, g, h \in R: (f \cdot g) \circ h = (f \circ h) \cdot (g \circ h)
\]

\vspace{5pt}

\noindent Let $d \colon R \to R$ be a function satisfying the chain rule, that is

\[
\forall f, g \in R: d(f \circ g) = \big(d(f) \circ g \big) \cdot d(g)
\]

\vspace{5pt}

\noindent Let $\int \colon R \to R$ be a function opposite to $d$. Then

\vspace{15pt}
\noindent\textbf{Theorem.} $\forall f, g, h \in R:$

\[
\big(\int f \cdot d(g)\big) \circ h = \int (f \circ h) \cdot d(g \circ h)
\]

\vspace{5pt}

\begin{lemmastar} $\forall f, g \in R:$ 

\[
\quad \int(f) \circ \int(g) = \int \big(f \circ \int(g) \cdot g \big).
\]

\vspace{10pt}
\textbf{Proof.}  For any $x, y \in R$, chain rule gives:

\[
d(x \circ y) = d(x) \circ y \cdot d(y).
\]  

\vspace{5pt}
Substitute $x = \int f$, $y = \int g$:  

\[
d\big(\int f \circ \int g\big) = \big(f \circ \int g \big) \cdot g.
\]  

\newpage
Applying $\int$ to both sides yields  

\[
\int f \circ \int g = \int \big(f \circ \int g \cdot g \big). \quad \square
\] 

\end{lemmastar}

\vspace{15pt}
\noindent\textbf{Proof of Theorem:}

\[
\big(\int f \cdot d(g)\big) \circ h = \big(\int f \cdot d(g)\big) \circ \int d(h)
\]

\vspace{5pt}
Applying Lemma:

\[
= \int\big(\big[f \cdot d(g)\big] \circ \big(\int d(h)\big) \cdot d(h)\big) 
\]  

\[
= \int\big(\big[f \cdot d(g)\big] \circ h \cdot d(h)\big)
\]  

\vspace{5pt}
By right distributivity in $R$:  

\[
= \int \big(\big[f \circ h \big] \cdot \big[d(g) \circ h \big] \cdot d(h)\big)
\]  

\vspace{5pt}
And reduce the expression by the chain rule for d:

\[
= \int \big(\big[f \circ h \big] \cdot d(g \circ h))
\]  

\vspace{5pt}
Or in other words:

\[
(f \triangleright g) \circ h = (f \circ h) \triangleright (g \circ h) 
\]

\vspace{5pt}
\centerline{ --- We call it "integral axiom". $\blacksquare$}

\vspace{10pt}
\noindent There is also a converse statement: informally speaking, for any distributive operation, 
there exists a corresponding function that satisfies the chain rule. 
For details see Theorem 2 in "On Algebraic Integral" \cite{umerov2025}.

\newpage
\section*{Examples}

We will now integrate several elementary functions using the integral axiom, 
to demonstrate how the theorem simplifies integration by turning it into a purely algebraic process.

\begin{example}
\[
\frac{1}{x \ln(x)} \triangleright x = \frac{1}{\ln(x)} \triangleright \ln(x)
\]

- put $\frac{1}{x}$ under the differential

\[
= [\frac{1}{x} \circ \ln(x)] \triangleright \ln(x) = (\frac{1}{x}  \triangleright x) \circ \ln(x)
\]

- apply the integral axiom.

\[
= \ln(\ln(x))
\]

- for details about integrating of $\frac{1}{x}$ see Lemma 4.8 in "On Algebraic Integral" \cite{umerov2025}


\end{example}

\begin{example}

\[
x^2 (x^3 + 1)^\frac{1}{3} \triangleright x = (x^3 + 1)^\frac{1}{3} \triangleright \frac{1}{3} x^3
\]

- put $x^2$ under the differential

\[
= \frac{1}{3}[x^\frac{1}{3} \circ (x^3 + 1) \triangleright (x^3 + 1)] = \frac{1}{3}[(x^\frac{1}{3} \triangleright x) \circ (x^3 + 1)]
\]

- use homogeneity of $d$ and the axiom of integral.  Note that composition has priority over the integral.

\[
= \frac{1}{3}[(\frac{3}{4} x^\frac{4}{3}) \circ (x^3 + 1)] = \frac{1}{4} (x^3 + 1)^\frac{4}{3}
\]

\end{example}

\begin{example}

\[
\frac{\sin(x)}{1 + \cos^2(x)} \triangleright x = \frac{1}{1 + \cos^2(x)} \triangleright (-\cos(x))
\]

\[
= -(\frac{1}{1 + x^2} \circ \cos(x) \triangleright \cos(x)) = -(\frac{1}{1 + x^2} \triangleright x) \circ \cos(x)
\]

- the integral axiom

\[
= -\arctan(\cos(x))
\]

\end{example} 


Simplicity of finding integrals demonstrated in Examples 1–3 suggests that the 
proposed algebraic method of integration can be turned into a fast and efficient symbolic 
integration algorithm for computer algebra systems and symbolic libraries such as MATLAB and SymPy. 
It remains an open question whether the method can accelerate numerical integration — which, in theory, 
could lead to the optimization and acceleration of neural networks and other mathematical systems.


\begin{thebibliography}{99}

\bibitem{umerov2025}
N. Umerov
\textit{On Algebraic Integral}.
Zenodo, 2025.
DOI: 10.5281/zenodo.17376700

\end{thebibliography}

\end{document}